\documentclass[12pt]{article}

\usepackage{sbc-template}

\usepackage[brazilian]{babel}
\usepackage[T1]{fontenc} % Codificação para fontes com suporte a caracteres acentuados
\usepackage[utf8]{inputenc} % Codificação UTF-8 (necessário apenas para pdflatex)
\usepackage{multicol}
\usepackage{multirow}
\usepackage{tabularx}
\usepackage{graphicx,url}
\usepackage{float}
\usepackage{amsmath,amssymb,amsfonts}
\usepackage{algorithmic}
\usepackage{textcomp}
\usepackage{enumerate}
\usepackage{soul}
\usepackage{verbatim}
\usepackage[dvipsnames,table,xcdraw]{xcolor}
\usepackage[shortlabels]{enumitem}
\usepackage{scalefnt} % pacote para diminuir fonte das tabelas
\usepackage{comment} % pacote de comentarios
\usepackage{lscape}
\usepackage{fancyhdr}
\usepackage{sbgames}
\usepackage[super]{nth}

% hyperref tem de vir por último, depois de todos os outros pacotes.
\usepackage{hyperref}
% DEBUG: colore citações, remissões internas e URLs. Comente a linha abaixo
% e descomente a seguinte para gerar a versão de submissão, em preto.
\hypersetup{colorlinks=true, citecolor=OliveGreen, linkcolor=BrickRed, urlcolor=MidnightBlue}
%\hypersetup{hidelinks}

%previne hifenização
\tolerance=1
\emergencystretch=\maxdimen
\hyphenpenalty=10000
\hbadness=10000

\sloppy


\trilha{Cultura} %Preencha com: Artes \& Design | Computação | Cultura | Saúde | Indústria | Educação | Tutoriais | MAGICA | WGP-Jogos
\local{Rio de Janeiro/RJ} % Exemplo: Manaus/AM
\ano{2026} % Exemplo: 2024
\edicao{Edição XXV} %Exemplo: XXIII

\title{Miaugerdon: desinformação, alteridade e construção do preconceito em um serious game \\
\medskip
\small{
    \textit{Title: Miaugerdon: misinformation, otherness, and the construction of prejudice in a serious game}
}
}

\author{Isadora de Moura Tebaldi\inst{1}, Jonatan de Almeida Torres Souza\inst{1}}

\address{Engenharia de Sistemas e Computação PESC/Coppe\\
  Universidade Federal do Rio de Janeiro -- Rio de Janeiro -- Brasil
  \email{itebaldi@cos.ufrj.br, jonatantorres@cos.ufrj.br}
}


\begin{document}

\maketitle

\thispagestyle{plain}

\begin{abstract}
\textbf{Introduction:} Perceptions about someone little known can be shaped by the information received about them, and misinformation may involve not only false content but also generalization, repetition and the delegitimization of sources.
\textbf{Objective:} This paper presents the design of Miaugerdon, a serious game that investigates how narrative and mechanics represent the formation of prejudice mediated by misinformation.
\textbf{Methodology or Steps:} In the game, a cat forms a negative perception of another through the discourse of a source he trusts. The work maps the game elements to their functions and analyses the design choices.
\textbf{Expected Results:} The design articulates misinformation as a gradual process, the construction of the other before direct contact and the player's agency in the revision of beliefs. Without a systematic evaluation with players, planned as future work, no effect on attitudes is claimed.
\end{abstract}

\keywords{Serious games, Misinformation, Prejudice, Otherness, Game design.}

\begin{resumo}
\textbf{Introdução:} A percepção sobre alguém pouco conhecido pode ser moldada pelas informações recebidas sobre ele, e a desinformação pode envolver não só conteúdos falsos, mas generalização, repetição e deslegitimação de fontes.
\textbf{Objetivo:} Este artigo apresenta a concepção de Miaugerdon, um serious game que investiga como narrativa e mecânicas representam a formação de preconceitos mediada pela desinformação.
\textbf{Metodologia ou Etapas:} No jogo, um gato forma uma percepção negativa sobre outro pelo discurso de uma fonte em que confia. O trabalho mapeia os elementos do jogo às suas funções e analisa as escolhas de design.
\textbf{Resultados Esperados:} O design articula a desinformação como processo gradual, a construção do outro antes do contato direto e a agência do jogador na revisão de crenças. Sem avaliação sistemática com jogadores, prevista como trabalho futuro, nenhum efeito sobre atitudes é afirmado.
\end{resumo}

\palavraschave{Serious games, Desinformação, Preconceito, Alteridade, Design de jogos.}

\section{Introdução}

A percepção sobre pessoas ou grupos pouco conhecidos pode ser moldada pelas informações recebidas sobre eles. Generalizações, estereótipos e discursos de ameaça podem produzir julgamentos antes mesmo do contato direto. Nesse processo, a desinformação pode envolver não apenas conteúdos falsos, mas também distorção, repetição, generalização e deslegitimação de fontes \cite{bad_news}.

Jogos digitais permitem explorar essas dinâmicas de forma interativa, colocando o jogador diante de informações, relações e consequências. Serious games têm sido utilizados para promover reflexão social \cite{dallaqua}, empatia e tomada de perspectiva \cite{Belman2009}, além de abordar preconceitos contra grupos percebidos como diferentes \cite{a_path_out}. Estudos sobre contato intergrupal também indicam que a interação entre grupos pode contribuir para reduzir atitudes preconceituosas \cite{Pettigrew2006}.

Partindo dessa perspectiva, este trabalho apresenta Miaugerdon, um jogo digital em que Caju, um gato que confia em Mr. T. como principal fonte de informação, passa a formar uma percepção negativa sobre um novo gato a partir de comentários e generalizações feitas por ele. Ao longo da experiência, informações e acontecimentos contraditórios levam Caju e o jogador a confrontar diferentes versões sobre o recém-chegado.

A proposta é representar mecanismos de influência, controle informacional e construção de preconceitos contra o diferente por meio da narrativa e das ações do jogador. Assim, o trabalho discute a concepção de \textit{Miaugerdon} como um \textit{serious game} que articula desinformação e preconceito, com foco em seu desenvolvimento e em suas escolhas de design. A partir disso, busca-se responder à seguinte pergunta de pesquisa: \textit{como narrativa e mecânicas de um serious game podem ser articuladas para representar a formação de preconceitos mediada pela desinformação?}

O artigo está organizado da seguinte forma: a Seção 2 apresenta a fundamentação teórica e os trabalhos relacionados; a Seção 3 descreve a concepção e reformulação de \textit{Miaugerdon}; a Seção 4 discute suas escolhas narrativas e mecânicas e apresenta as limitações do trabalho; e a Seção 5 reúne as conclusões e os trabalhos futuros.

\section{Fundamentação Teórica e Trabalhos Relacionados}

Esta seção apresenta os principais conceitos que fundamentam a proposta de Miaugerdon e discute trabalhos relacionados ao uso de jogos para reflexão social. São abordados estudos sobre serious games, empatia, preconceito e desinformação, buscando situar o jogo em relação às abordagens já exploradas na literatura.

\subsection{Serious games e reflexão social}
Os serious games são jogos desenvolvidos com objetivos que ultrapassam o entretenimento, podendo ser utilizados em contextos de aprendizagem, conscientização e transformação social. A revisão realizada por Dallaqua et al. mostra que esse campo reúne trabalhos voltados ao uso de jogos para abordar problemas sociais, explorar valores e estimular reflexão por meio de experiências interativas \cite{dallaqua}. Nesses jogos, a compreensão do tema pode surgir não apenas da apresentação de informações, mas também das situações, escolhas e consequências vivenciadas pelo jogador.

Belman e Flanagan discutem especificamente o uso de jogos para promover empatia e tomada de perspectiva, analisando exemplos como PeaceMaker, Hush e Layoff \cite{Belman2009}. Os autores destacam que a capacidade de um jogo de provocar reflexão depende de como narrativa, personagens e ações do jogador são articulados à experiência proposta. Dessa forma, o jogador pode ser colocado diante de perspectivas e conflitos distintos dos seus, participando ativamente das situações representadas.

\subsection{Preconceito, alteridade e contato com o diferente}
O preconceito pode se formar a partir de estereótipos, generalizações e percepções negativas atribuídas a indivíduos identificados como pertencentes a grupos externos ou diferentes. Quando há pouco contato direto entre as partes, essas percepções podem ser construídas principalmente a partir de informações recebidas de terceiros, reforçando interpretações prévias sobre o “outro”. Nesse sentido, Pettigrew e Tropp demonstram, por meio de uma ampla meta-análise, que o contato intergrupal está associado à redução de atitudes preconceituosas, especialmente por permitir que expectativas e imagens previamente estabelecidas sejam confrontadas pela experiência direta \cite{Pettigrew2006}.

Essa relação também tem sido explorada em jogos digitais. Cross et al. investigaram uma experiência em que jogadores assumiam a perspectiva de um refugiado sírio e observaram redução de atitudes preconceituosas e aumento de empatia após a interação \cite{a_path_out}. O estudo evidencia como jogos podem explorar atitudes dirigidas a grupos percebidos como externos, incluindo contextos relacionados à migração e ao preconceito contra refugiados.

Esses trabalhos reforçam a importância do contato e da tomada de perspectiva na revisão de julgamentos sobre grupos percebidos como diferentes, além de apontarem o potencial dos jogos como espaços para representar e explorar essas relações de forma interativa.

\subsection{Jogos e desinformação}
Jogos digitais também têm sido explorados como ferramentas para desenvolver resistência à desinformação. Roozenbeek e van der Linden apresentam o Bad News, jogo no qual o jogador assume o papel de um produtor de notícias falsas e utiliza estratégias como polarização, apelo emocional, teorias conspiratórias e descrédito de fontes \cite{bad_news}. Ao colocar o jogador em contato direto com essas técnicas, o jogo busca tornar mais reconhecíveis os mecanismos usados para manipular informações. Os resultados do estudo indicaram melhora na capacidade dos participantes de identificar conteúdos e estratégias enganosas, reforçando o potencial de abordagens interativas para trabalhar pensamento crítico e alfabetização midiática.

Enquanto alguns trabalhos abordam desinformação e outros investigam preconceito e empatia, \textit{Miaugerdon} articula esses temas ao representar como uma percepção negativa sobre o outro pode ser construída antes do contato direto, a partir de informações e generalizações de uma fonte manipuladora. A experiência explora, assim, a relação entre controle informacional, formação de estereótipos e revisão de julgamentos.

\section{Concepção do Miaugerdon}

Esta seção apresenta a concepção de \textit{Miaugerdon}, descrevendo sua premissa, personagens, estrutura narrativa e principais mecânicas. O foco está em como esses elementos foram organizados para representar processos de influência, desinformação e formação de percepções sobre o diferente.

\subsection{Visão geral e personagens}

O jogo acompanha Caju, um gato que vive com seu tutor humano, Alfredo, e descobre que ele pretende adotar Soneca, um novo gato vindo de um abrigo. Incomodado com a chegada do desconhecido, Caju recorre a Mr. T., um gato em quem confia e que passa a influenciar sua percepção sobre o recém-chegado por meio de comentários, generalizações e informações negativas. Mr. T. atua como principal agente de manipulação, enquanto Alfredo apresenta informações e situações que desafiam essa narrativa. Soneca, por sua vez, ocupa o papel daquele percebido como ``de fora'' ou diferente, sendo inicialmente conhecido por Caju por meio do discurso de terceiros.

Sob a influência de Mr. T., Caju inicia um plano para construir uma máquina de dominação mundial, que funciona como objetivo central da partida. A experiência é apresentada em 2D, com vista isométrica, e se passa principalmente na casa de Alfredo, explorada pelo jogador enquanto cumpre uma sequência de tarefas. Durante a partida, Alfredo circula pelo ambiente e pode interromper ações consideradas suspeitas, levando o jogador a alternar entre atividades relacionadas ao plano e comportamentos comuns de um gato. A progressão combina exploração, interação com objetos e personagens, distrações e um sistema de suspeita.

Apesar dos temas abordados, a experiência adota um tom predominantemente cômico. A ideia de um gato construindo uma máquina de dominação mundial, combinada a comportamentos cotidianos utilizados para disfarçar suas ações, introduz situações exageradas e absurdas ao longo da partida. O humor funciona como parte da identidade do jogo e como recurso narrativo para abordar os temas propostos sem abandonar seu caráter lúdico.

\subsection{Estrutura narrativa}

A narrativa se inicia quando Caju conta a Mr. T. que Alfredo pretende adotar um novo gato. A partir dessa informação, Mr. T. começa a apresentar comentários, generalizações e notícias falsas sobre o recém-chegado, construindo uma percepção negativa antes que Caju tenha qualquer contato direto com ele. Falas como ``esses gatos são primitivos'' e ``eles vão roubar sua cama'' reforçam a ideia de ameaça e contribuem para que Caju passe a enxergar a chegada de Soneca de forma negativa. Sob essa influência, o protagonista decide construir uma máquina de dominação mundial.

A progressão ocorre por meio de pequenas tarefas necessárias para a construção da máquina. Ao longo dessas etapas, Caju retorna periodicamente a Mr. T. para relatar seu progresso e receber novas orientações, reforçando a relação de confiança e dependência entre os personagens. Paralelamente, o protagonista passa a encontrar informações, falas e acontecimentos que contradizem o discurso de Mr. T. Diante dessas contradições, o personagem também busca deslegitimar fontes alternativas, por meio de afirmações como ``essa mídia é comprada'' e ``você não deveria se informar por ela'', ampliando a manipulação para além da circulação de informações falsas.

À medida que essas contradições se acumulam, Caju começa a reconsiderar as informações recebidas e a influência exercida por Mr. T. Caso avance até as etapas finais do plano, o protagonista revisita os acontecimentos vivenciados e se depara com a decisão de concluir ou abandonar a construção da máquina.

A experiência foi concebida para partidas curtas, com duração aproximada de cinco minutos, e apresenta quatro desfechos possíveis. Caju pode concluir a construção da máquina, desistir do plano, ser descoberto por Alfredo enquanto executa as tarefas ou não conseguir concluí-las antes do término do tempo disponível. Os finais decorrem tanto das decisões do jogador quanto de seu desempenho ao longo da partida.

\subsection{Mecânicas e objetivos de design}
\begin{table}[ht]
\centering
\caption{Elementos do jogo e suas funções na experiência.}
\label{tab:elementos_jogo}
\begin{tabularx}{\linewidth}{|X|X|}
\hline
\textbf{Elemento do jogo} & \textbf{Função na experiência} \\ \hline
Repetição de tarefas para Mr. T. & Construção de confiança na personagem \\ \hline
Informações negativas sobre o novo gato & Formação de percepção sem contato direto \\ \hline
Generalizações sobre ``esses gatos'' & Representação de estereótipos \\ \hline
Notícias contraditórias & Exposição a versões alternativas \\ \hline
Deslegitimação da imprensa & Controle das fontes de informação \\ \hline
Interações com Alfredo & Contato com uma perspectiva diferente \\ \hline
Lembranças durante a montagem & Revisão das crenças anteriores \\ \hline
\end{tabularx}
\end{table}

A progressão do jogador é condicionada por um sistema de suspeita associado a Alfredo. Enquanto Caju executa ações relacionadas ao plano, a barra de suspeita aumenta e pode ser reduzida por meio de comportamentos comuns de um gato. Alfredo circula pela casa seguindo trajetos de forma aleatória, mas seu comportamento se altera conforme o nível de suspeita. Quando a barra atinge o nível amarelo, ele se desloca até Caju para verificar suas ações; no nível vermelho, passa a segui-lo pela casa. Caso Caju seja flagrado executando uma ação relacionada ao plano, a partida pode ser encerrada.

Essa mecânica exige que o jogador alterne entre avançar na construção da máquina e realizar ações cotidianas para evitar a atenção de Alfredo. Além de introduzir uma condição de risco à progressão, esse contraste entre um plano de dominação mundial e comportamentos típicos de um gato contribui para o tom cômico da experiência.

As demais mecânicas e situações narrativas foram organizadas de modo a relacionar as ações do jogador aos temas centrais do jogo. A repetição de tarefas para Mr. T., por exemplo, busca construir gradualmente uma relação de confiança, enquanto as informações negativas sobre Soneca surgem antes de qualquer contato direto com o personagem. Generalizações, notícias contraditórias e a deslegitimação de fontes representam diferentes formas de construção e manutenção de uma percepção negativa sobre o outro.

Em contraponto, as interações com Alfredo introduzem perspectivas alternativas, enquanto as lembranças retomadas durante a montagem da máquina permitem que Caju confronte informações recebidas anteriormente. A Tabela~\ref{tab:elementos_jogo} sintetiza a relação entre esses elementos e suas funções na experiência.

\section{Discussão}

Esta seção discute como as escolhas narrativas e mecânicas do jogo representam os processos de desinformação, construção de preconceitos e revisão de percepções apresentados ao longo da experiência. A análise considera especialmente o papel de Mr. T. como fonte manipuladora, a construção de Soneca como figura do ``outro'' e a forma como as ações do jogador acompanham a transformação das crenças de Caju.

\subsection{Desinformação como processo, não como notícia falsa}
No jogo, a desinformação não é representada apenas pela circulação de informações falsas, mas por um conjunto de estratégias que moldam a forma como Caju interpreta os acontecimentos. Mr. T. utiliza generalizações, repetição de discursos negativos e deslegitimação de fontes alternativas, aproximando-se de técnicas de manipulação discutidas em trabalhos como o \textit{Bad News} \cite{bad_news}. Dessa forma, o problema deixa de estar restrito à veracidade de uma informação isolada e passa a envolver também quais fontes são consideradas confiáveis e quais interpretações são reforçadas.

Essa escolha permite representar a desinformação como um processo gradual. A influência de Mr. T. é construída ao longo das interações e das sucessivas atualizações fornecidas por Caju, até que informações contraditórias começam a enfraquecer essa relação de confiança. A progressão narrativa busca, assim, mostrar como percepções podem ser formadas e mantidas por um ambiente informacional controlado antes de serem confrontadas por outras experiências.

\subsection{Construção do ``outro''}

Soneca ocupa na narrativa a posição daquele percebido como ``de fora''. Antes de qualquer contato direto, Caju forma uma imagem sobre o recém-chegado a partir das generalizações e informações negativas apresentadas por Mr. T. Dessa forma, o preconceito é construído não pela experiência com Soneca, mas por um discurso que o associa previamente a características negativas e à ideia de ameaça.

Embora Soneca não represente diretamente um grupo social específico, essa dinâmica pode remeter a diferentes formas de exclusão social, entre elas a xenofobia, nas quais indivíduos percebidos como externos ao grupo são julgados a partir de estereótipos antes do contato direto. A opção por não vincular o personagem a um grupo específico preserva a alegoria do ``outro'' e permite que a experiência seja interpretada a partir de diferentes contextos de alteridade.

Essa escolha narrativa dialoga com estudos sobre contato intergrupal, que associam a interação entre grupos à revisão de percepções e atitudes preconceituosas \cite{Pettigrew2006}, bem como com experiências em jogos que exploram preconceito contra grupos percebidos como externos \cite{a_path_out}. Em \textit{Miaugerdon}, esse processo aparece no contraste entre aquilo que Caju ouve sobre Soneca e as informações e experiências que gradualmente desafiam essa imagem inicial.

\subsection{Envolvimento do jogador e reflexão}

A proposta também busca envolver o jogador no próprio processo de influência vivido por Caju. Em vez de apenas observar a manipulação exercida por Mr. T., o jogador executa as tarefas sugeridas por ele, acompanha o avanço do plano e retorna periodicamente para relatar seu progresso. Dessa forma, a agência do jogador é incorporada à construção da confiança e à continuidade das ações de Caju.

À medida que surgem informações e experiências contraditórias, essa mesma agência passa a acompanhar a revisão das crenças do protagonista. A montagem da máquina funciona como ponto de convergência desse processo, pois retoma decisões e acontecimentos anteriores e coloca em questão a continuidade do plano. Assim, a experiência busca fazer com que a reflexão não ocorra apenas no nível narrativo, mas também a partir das ações que o próprio jogador realizou ao longo da partida.

A experiência também adota um tom cômico, presente no caráter exagerado do plano de Caju e na própria construção de uma máquina de dominação mundial. O uso do humor não elimina o propósito reflexivo da experiência, mas constitui uma escolha de design para mediar temas sociais por meio de uma narrativa lúdica e exagerada.

\subsection{Limitações}

O protótipo não foi submetido a uma avaliação sistemática com jogadores, portanto o presente trabalho não permite afirmar que a experiência modifica atitudes, reduz preconceitos ou aumenta a resistência à desinformação. A proposta deve ser entendida como uma experiência projetada para estimular reflexão sobre esses temas, e não como evidência de mudança comportamental ou atitudinal.

\section{Conclusão e Trabalhos Futuros}
Este trabalho apresentou e analisou o design de um \textit{serious game} que articula desinformação e construção do preconceito contra o diferente por meio de narrativa e mecânicas interativas. A proposta utiliza a relação entre Caju e Mr. T., a construção da imagem de Soneca e a revisão gradual das crenças do protagonista para representar como percepções negativas podem ser formadas e posteriormente questionadas.

Como trabalho futuro, pretende-se conduzir uma avaliação sistemática com jogadores, combinando observação e entrevistas após a experiência. Essa etapa poderá investigar se os participantes perceberam a manipulação exercida por Mr. T., como interpretaram Soneca, em que momento passaram a desconfiar das informações recebidas e de que forma compreenderam a alegoria proposta. Também será possível comparar percepções antes e depois dos principais acontecimentos narrativos e analisar diferenças de interpretação entre os jogadores.

\section*{Declaração de Uso de Inteligência Artificial}
Em conformidade com as práticas de transparência e integridade acadêmica, declara-se que modelos de linguagem de grande porte (LLMs) foram utilizados como ferramenta de apoio neste trabalho.

Especificamente, essas ferramentas auxiliaram na redação e revisão do texto, na organização da estrutura do artigo e no refinamento da concepção narrativa e das decisões de design relacionadas à reformulação do protótipo.

Todo o conteúdo produzido com auxílio de IA foi revisado, validado e, quando necessário, corrigido pelos autores, que assumem integral responsabilidade pela concepção, pelo conteúdo e pelas conclusões apresentadas.

\bibliographystyle{sbc}
\bibliography{refs}

\end{document}
